% =====================================================================
% Scrubbing the Run: Replay-Native Provenance Visualization for
% Closed-Loop Beamline Experiments
%
% Target venue: VAxAutoSci (Visual Analytics in the Age of Autonomous
% Science), an IEEE VIS 2026 workshop. Short paper, 4-6 pages excluding
% references, two-column VGTC conference style.
%
% This is the submission file. The prose lives in sections/*.tex so that
% this file and the generic-layout preview.tex read the SAME text.
%
% TEMPLATE: this file assumes the official IEEE VIS / VGTC conference
% LaTeX template. Download the VIS 2026 author kit (vgtc.cls, the
% abbrv-doi .bst, and supporting .sty files), place them in this folder,
% and reconcile the front-matter macros below with that kit. We do not
% vendor the class files. Until they are present this file will not
% compile; use preview.tex for a quick read (see README.md).
% =====================================================================

\documentclass[journal]{vgtc}                % per the VIS author kit

\usepackage{mathptmx}
\usepackage{graphicx}
\usepackage{times}
\usepackage{booktabs}
\usepackage{url}
\usepackage{microtype}

% --- VGTC front matter (reconcile names with the author kit) -----------
\onlineid{0}
\vgtccategory{Research}
\vgtcpapertype{application/design study}
\authorfooter{
  \item Author affiliation and contact. Placeholder for submission.
}

\title{Scrubbing the Run: Replay-Native Provenance Visualization for
Closed-Loop Beamline Experiments}

\author{Anonymous or named per single/double-blind choice}

\abstract{Autonomous and closed-loop scientific instruments, from self-driving labs to
conducted synchrotron experiments, generate provenance at machine speed: dense
in decisions, multi-pass, and largely invisible to the humans who must later
audit it. A scientist returning to a finished run asks plain questions: why did
it stop here, which pass converged, can I prove what it actually did?
Conventional post-hoc provenance graphs flatten the temporal and causal texture
those questions depend on. We observe that event-sourced systems of record
already answer such questions internally, because they reconstruct any past
state by folding an immutable, ordered event stream. We turn that fold into an
interaction. This paper contributes (1) a task abstraction for auditing
closed-loop runs, and (2) a \emph{replay scrubber}: a draggable time cursor
that replays a run to any recorded instant, showing per-step activity
swim-lanes, nested convergence brackets colored by verdict, and a
content-addressed fidelity badge that lets the view verify the replayed state
matches the record. We ground the prototype in a real closed-loop autofocus
alignment on the APS 2-BM micro-CT beamline, and discuss how event-sourced
records reshape provenance visualization for autonomous science.
}

% \CCScatlist{ ... }   % add per the author kit
% \teaser{ ... }       % optional teaser figure

\begin{document}

\maketitle

\firstsection{Introduction}
\label{sec:intro}
% Introduction body. The section heading is issued by the driver file
% (main.tex uses \firstsection, preview.tex uses \section).

A modern beamline increasingly runs itself. Alignment, focusing, and in some
settings the choice of what to measure next are handed to software that
decides and acts at machine speed, whether in a self-driving materials
laboratory or in a synchrotron procedure conducted on the scientist's behalf.
When such a run ends, or stops unexpectedly, a human still has to answer for
it. The first questions are plain: why did it stop here, which pass finally
converged, and can I prove what the instrument actually did?

These are provenance questions, but the provenance of an autonomous run is
unlike the human-authored record most tools were built for. It arrives at
machine speed, it is dense in decisions, and it is multi-pass: a single
alignment may bracket a peak over several iterations, each with its own
verdict. Post-hoc provenance graphs and flat logs flatten exactly the
temporal and causal texture these questions depend on. When a run is
interrupted the log simply stops, leaving the auditor to guess what was in
flight and whether it took effect, and a static record offers no way to show
that it has not been altered after the fact.

We observe that one class of system already answers such questions
internally. An event-sourced system of record never stores current state
directly; it reconstructs any past state on demand by folding an immutable,
ordered stream of events. Our central idea is to expose that fold as an
interaction. The result is a \emph{replay scrubber}: a draggable time cursor
that an auditor moves through a recorded run, with the system folding the
event stream to the chosen instant and rendering the state as it was then.
Because the same fold that makes the backend correct is the one driving the
view, replay is faithful by construction, and a content-addressed check of
the recorded expansion lets the scrubber verify, rather than merely assert,
that the state shown at the cursor matches the record.

We ground the design in a real run from an operating instrument: a
CORA-conducted autofocus alignment on the APS 2-BM micro-CT beamline, a
four-iteration peak-bracket search with per-step activities and a converged
verdict on the final pass. We are deliberate about scope. This is a
\emph{conducted}, closed-loop procedure; the fully agent-closed
perceive-decide-act loop, in which an AI selects each next action, is out of
scope, and the prototype is a research visualization over exported run data
rather than a deployed interface.

\medskip\noindent This paper makes four contributions:
\begin{itemize}
  \item \textbf{Replay-native provenance (insight).} We reframe the
        fold-on-read mechanism of an event-sourced system of record as an
        interaction primitive for provenance: the fold the backend uses to
        reconstruct state becomes a control the auditor drags through
        recorded time, rather than a static post-hoc graph to read.
  \item \textbf{A task abstraction for auditing closed-loop runs.} We
        characterize four recurring audit tasks and map each to an
        event-log query and a visual encoding.
  \item \textbf{The replay scrubber (technique).} We design a single
        timeline that unifies per-step activity swim-lanes, nested
        convergence brackets colored by verdict, open-interval markers for
        unfinished steps, a draggable fold-to-version cursor, and a
        content-addressed fidelity badge.
  \item \textbf{A grounded case study.} We demonstrate the tasks on a real
        closed-loop autofocus alignment on the APS 2-BM micro-CT beamline,
        including an interrupted-run scenario in which the scrubber localizes
        the exact unfinished step.
\end{itemize}


\section{The Auditor's Tasks}
\label{sec:tasks}
% The Auditor's Tasks body.

We derive four recurring tasks, which we label T1 through T4, from what a
beamline scientist needs when returning to a conducted run, and make each
concrete as a query over the event log paired with a visual encoding
(Table~\ref{tab:tasks}). The tasks are
not hypothetical: each is a question we ask of the lights-out, agent-supervised
run in Section~\ref{sec:scrubber}, and each is answerable only because the record is
an ordered, immutable event stream rather than a final-state snapshot.

\textbf{T1, recover the interrupted state.}
\emph{What was in flight when the beam dropped?} The system records intent
before effect, so an unfinished step appears in the log as an activity that
was opened but never closed. The query finds the activity with no matching
outcome; the encoding draws it as an open interval, a visible gap rather than
a silent truncation.

\textbf{T2, reproduce a decision.}
\emph{Replay the run to the instant a choice was made, and show the state
then.} The query folds the event stream up to the chosen position; the
encoding is the draggable cursor together with the reconstructed state read
out beside it.

\textbf{T3, confirm convergence.}
\emph{Did this loop converge, and on which pass?} The query reads the
per-iteration verdict carried by each iteration-boundary event; the encoding
colors nested brackets by that verdict, so the bracketing search is legible at
a glance.

\textbf{T4, confirm the record is intact.}
\emph{Can I trust that nothing was altered?} The query recomputes the
content-addressed hash recorded with the run and compares it at the cursor;
the encoding is an inline fidelity badge that turns ``trust me'' into ``the
view checked it.''

These four tasks order the rest of the paper: T2 and T3 drive the scrubber
itself (Section~\ref{sec:scrubber}), T1 drives the interruption-recovery
walkthrough (Section~\ref{sec:crash}), and T4 rests on the immutable substrate
(Section~\ref{sec:substrate}).

\begin{table}[t]
  \centering
  \caption{Four auditor tasks, each mapped to an event-log query and a visual
  encoding. The mapping is the spine of the design.}
  \label{tab:tasks}
  \footnotesize
  \begin{tabular}{@{}p{0.235\linewidth}p{0.33\linewidth}p{0.33\linewidth}@{}}
    \toprule
    \textbf{Task} & \textbf{Event-log query} & \textbf{Visual encoding} \\
    \midrule
    T1 Recover interrupted state &
      Activity opened with no closing outcome &
      Open interval (a visible gap) \\[0.4em]
    T2 Reproduce a decision &
      Fold the stream to position $t$ &
      Draggable cursor + folded state \\[0.4em]
    T3 Confirm convergence &
      Per-iteration verdict on boundary events &
      Nested brackets colored by verdict \\[0.4em]
    T4 Confirm integrity &
      Recompute and compare the recorded hash &
      Inline fidelity badge \\
    \bottomrule
  \end{tabular}
\end{table}


\section{The Replay Scrubber}
\label{sec:scrubber}
\textit{[Draft pending. Figure~1 (centerpiece, real exported data) plus a
Figure~2 inset (the fidelity badge): a time axis with per-step activity
swim-lanes, nested iteration brackets colored by the converged verdict,
in-flight steps drawn as open intervals, and a draggable fold-to-version
cursor whose position selects the folded state shown alongside. Walk through
the 2-BM autofocus run, answering T2 (reproduce a decision) and T3 (confirm
convergence).]}


\section{Crash Recovery}
\label{sec:crash}
% First pass written without the exported figure data; Figure 3 pending.
% TODO(data): name the exact step index, activity kind, and timestamp of the
% dangling in-flight record from the simulated-crash export.

The most operationally pressing question an auditor asks is also the one
conventional records answer worst: when a run stops unexpectedly, what was in
flight, and did it take effect? We reproduce that situation directly. Starting
from the recorded autofocus run, we drop the closing outcome of one activity,
simulating a conductor that died after recording its intent to act but before
recording the result, and replay the truncated stream through the scrubber
(Figure~3).

Because the conductor writes intent before effect, the unfinished activity is
still present in the log, just without its closing event. The fold therefore
reconstructs a state in which exactly one activity is open, and the scrubber
draws it as an open interval at its lane and time, with the convergence bracket
around it left unclosed. Rather than a timeline that simply ends, the auditor
sees precisely which step was mid-flight and where, the payoff for T1. Dragging
the cursor to the instant just before the gap shows the last state known to be
complete, which is the state a recovery would resume from. This is the case
post-hoc provenance graphs handle least gracefully, since a graph assembled
after the fact has no node for an action that never reported a result; the
intent-before-effect record, read by replay, makes the missing result itself
legible.


\section{Substrate}
\label{sec:substrate}
\textit{[Draft pending. Figure~5, the single architecture figure, labeled
``substrate, not contribution'': the event record behind the scrubber, with
just enough of the envelope, the run-start payload, and the decision record to
show where the fold and the recorded hash come from. Immutability is stated in
one sentence supporting the fidelity badge, with a citation to the append-only
store, not given its own section.]}


\section{Related Work}
\label{sec:related}
% Related Work body.

Our work draws on five lines of research, none of which composes its central
primitive: a deterministic fold over an immutable event record, exposed as a
scrub interaction, carrying a content-addressed fidelity check and provenance
semantics for autonomous loops.

\paragraph{Provenance visualization and notebook time travel.}
Retrospective provenance, the record of what an execution actually did, has
long been distinguished from the prospective recipe that specifies
it~\cite{lim2010}. Systems that visualize it, such as
VisTrails~\cite{vistrails2006} and Loops~\cite{loops2024}, let an analyst
navigate a version tree and compare or diff recorded versions of a workflow or
notebook. Verdant~\cite{verdant2018, verdant2019} is the closest prior
interaction to ours: it pairs recorded notebook history with a draggable
timeline cursor. Its replay, however, reconstructs a past version by
re-executing it and explicitly ignores the dependencies the older version
originally had, so it cannot guarantee that the reconstructed state matches
what happened, and it targets a single analyst's local notebook rather than an
autonomous run. QIIME~2 Provenance Replay~\cite{qiime2replay} likewise turns a
recorded provenance graph into regenerated code that is re-run offline, not an
interactive view. We differ in that we reconstruct state by folding an
immutable event record and certify the match at the cursor, rather than
navigating versions or replaying by re-execution.

\paragraph{Record-replay and time-travel debugging.}
Omniscient debuggers such as Pernosco and rr~\cite{rr2017}, the
Whyline~\cite{whyline2008}, and reducer-based state containers share our
record-once, replay-later posture and give instant access to any past instant.
In all of them, deterministic replay is held as an internal guarantee and is
never rendered to the user. In short, these systems replay computation, whereas
we replay recorded run state with provenance semantics and a per-cursor,
content-addressed integrity check that makes the faithfulness of the replay a
visible, checkable affordance.

\paragraph{Tamper-evident and verifiable provenance.}
The canonical integrity structures, Crosby and Wallach's history
tree~\cite{crosby2009} and Certificate Transparency's Merkle logs~\cite{ct2013},
verify membership and consistency as machine-checked proofs with no user
interface; content-addressed identifiers such as SWHID~\cite{swhid} expose the
same recompute-and-compare primitive only as a command-line property. The
closest human-facing analog is C2PA Content Credentials~\cite{c2pa}, which
renders the authenticity of recorded media; content-credential provenance has
recently been carried to machine-learning pipelines~\cite{atlas2025}, where a
content-addressed hash-versus-reference check runs at a programmatic service
layer over discrete artifacts, with no visualization and no notion of replayed
state. We differ in that we render a content-addressed verification verdict to
a human, at each scrubbed instant, over reconstructed run state.

\paragraph{Reproducibility records and badges.}
Provenance-record tools such as CPR~\cite{cpr2021} and Whole Tale~\cite{wholetale},
workflow-run provenance crates~\cite{wrroc2024}, and artifact
badges~\cite{acmbadges} either emit static reports and immutable output
snapshots, or certify a run through human audit with results required to agree
only within a tolerance. None reconstructs per-instant state or renders an
exact, content-addressed match. We differ in that our fidelity badge is an
automated equality check recomputed as the user scrubs, not a coarse,
evaluation-time certification.

\paragraph{Conformance checking.}
Process mining supplies the canonical quantitative verdict of whether a run
conformed, through token-replay fitness~\cite{rozinat2008} and localized
alignment diagnostics~\cite{berti2019, pm4py}; online variants detect
deviations on a live stream~\cite{vanzelst2019}. These verdicts are encoded as
static overlays on a process model, or as per-move type coloring along a
completed trace, never as verdicts along an iteration axis with interval
structure, and alignments computed post hoc on complete traces have no
representation for an interrupted step. Tellingly, recent surveys of
conformance visualization conclude that the visual encoding of conformance
results remains an open problem~\cite{rehse2022, hagerehse2025}. We contribute
exactly such an encoding: nested convergence brackets colored by per-iteration
verdict, with open intervals for in-flight steps.

\paragraph{Visual analytics for autonomous experimentation.}
Self-driving-lab and autonomous-beamline systems increasingly keep complete
operations logs and present live dashboards. VISION~\cite{vision2025} logs
operations for backend monitoring, while optimizer and campaign tools such as
Olympus~\cite{olympus2021}, the autonomous SAXS/WAXS framework at ALS and
PETRA~III~\cite{alspetra2026}, Tsuchinoko~\cite{tsuchinoko},
ChemOS~2.0~\cite{chemos2}, A-Lab~\cite{alab2023}, and Bluesky's live
callbacks~\cite{bluesky} render score and posterior fields, convergence lines,
and live plots. These are forward-streaming views of a run in progress; none
offers a scrubbable, post-hoc audit of the recorded run, per-iteration
convergence verdicts, or a fidelity check. That gap, and the audit and steering
tasks it implies, is what our replay scrubber addresses.


\section{Limitations and Future Work}
Our evaluation is a structured walkthrough of the four auditor tasks on one
recorded run, not a controlled study. We report a single run, on a single
instrument, and we have not yet run a user or expert study; whether the scrubber
speeds real audits, and which encodings practitioners prefer, are open empirical
questions. The interface is a research prototype over exported run data rather
than a deployed beamline tool, so we make no claims about throughput, scale, or
integration cost.

Two scope boundaries matter for how the contribution should be read. First, the
run is \emph{conducted} and \emph{agent-supervised}, not agent-closed: CORA
orchestrates the alignment and records every step, and a RunSupervisor agent
holds and resumes the run on beam conditions, but neither chooses the science.
That supervisor is a deterministic, rule-based agent making recorded decisions;
the figures show a genuine human-agent handoff, not a model selecting
experiments, and we are careful not to claim the latter. Second, the fidelity
badge attests \emph{faithfulness}, that the replayed state matches the recorded
expansion, by recomputing content-addressed hashes against an append-only store.
It is an integrity check, not a defense against an adversary with write access;
authenticity guarantees would need a separate trust anchor.

The fold-to-version helper that drives the cursor was written for this
visualization on top of the system's existing fold-on-read path. Because that
path keeps no snapshots, the cost of reconstructing a distant instant grows with
the length of the stream; for the short run shown here this is imperceptible, but
snapshotting for long campaigns is future work. Three further extensions are
natural and named, not built: a live now-edge that subscribes to the event tail
so the same timeline monitors a run in progress; run-to-run diffing so two
campaigns can be compared on one axis; and a deployed read interface, of which
the first read-only slice is on the roadmap. We also expect the design to
generalize beyond this run to other conducted procedures and other facilities,
since nothing in the scrubber is specific to it beyond the lifecycle, iteration,
and activity vocabulary it reads from the event log.


\bibliographystyle{abbrv-doi}                % per the VIS author kit
\bibliography{refs}

\end{document}
